\documentclass{article}
\usepackage[slovene]{babel}

\usepackage{ragged2e,lmodern} %sodobni LaTeX

\usepackage{array,booktabs, multirow} % Tabele
\usepackage{graphicx} % Slike

\usepackage{biblatex}
\addbibresource{literatura.bib} %bibliografije

\usepackage{hyperref} % Hiperpovezave




\usepackage{lipsum} % Lorem ipsum

\begin{document}
\tableofcontents
\section{Sklicevanje}
\label{sec:sklicevanje}

\section{Literatura}
Ta besedilo je del razdelka literatura. Seznam referenc bo natisnjen spodaj, vendar ker uporabljam možnost \texttt{heading=none}, naslov ni natisnjen.

\nocite{*}
% \printbibliography[title={Literatura, privzeto}]
% \printbibliography[title={Literatura, v kazalu}, heading=bibintoc]
% \printbibliography[title={Literatura, kot podrazdelek -- z zvezdico --}, heading=subbibliography]
% \printbibliography[title={Literatura kot razdelek, s številko}, heading=bibnumbered]
% \printbibliography[title={Literatura kot podrazdelek, brez št.-a, v kazalu}, heading=subbibintoc]
% \printbibliography[title={Literatura kot podrazdelek, s št.-o, v kazalu}, heading=subbibnumbered]
\printbibliography[title={Čeprav napišem naslov, se nič ne natisne}, heading=none]


% S uporabo ukaza \texttt{\textbackslash label\{sec:sklicevanje\}} smo označili ta razdelek, nanj pa se lahko sklicujemo z ukazom \texttt{\textbackslash ref\{sec:sklicevanje\}}. 

% Na primer, v razdelku \nameref{sec:sklicevanje} smo predstavili sklicevanje na različne elemente v dokumentu.
% Na primer, v razdelku \cref{sec:sklicevanje} smo predstavili sklicevanje na različne elemente v dokumentu.

\vspace{1cm}

% \begin{figure}
%   \centering
%   \includegraphics[width=3cm]{../02_latexOsnovnaOkolja/macka}
%   \caption{Slika čudovite mačke.}
%   \label{fig:macke}
%   \end{figure}
%   %...
%   Na sliki \ref{fig:macke} na strani \pageref{fig:macke} vidimo lepo mačko.



% \vspace{2cm}

% \begin{enumerate}
%   \item Oblecite nogavice.
%   \item \label{itm:cevlje} Obujte čevlje.
%   \item Pojdite na sprehod.
% \end{enumerate}
% %...
% V koraku \ref{itm:cevlje} ne poszabite zavezati vezalk.

% \vspace{2cm}

% Prva knjiga, ki sem jo prebral, je bila \cite{hp1}. 
% % ...
% \begin{thebibliography}{9}
%   \bibitem{hp1} J. K. Rowling, \textit{Harry Potter in kamen modrosti}, Založba EPTA, 1997.
%   \bibitem{hp2} J. K. Rowling, \textit{Harry Potter in dvorana skrivnosti}, Založba EPTA, 1998.
% \end{thebibliography}



\end{document}

